\section{The explicit offset correction and matrix constants}\label{ssm:sec:explicit}
\subsection{Why the finite saddle mass has the required rate}
Set
\[
 \theta_M=k-M/2+1,\qquad T_M=2\theta_M/\sqrt M.
\]
This is the exact parametrization $k=M/2-1+(T_M/2)\sqrt M$ used in \cite{mww}. Let $L_0(T)$ be the positive solution
\begin{equation}\label{ssm:eq:Lzero}
 \phi(L_0(T)+T)=L_0(T)\Phi(L_0(T)+T).
\end{equation}
In particular $L_0(0)=\ell$. In a neighborhood of zero, smoothness follows from the implicit function theorem because the derivative in $L$ does not vanish.

We carefully identify the finite saddle before importing a quantitative estimate. In \cite[Section 6]{mww}, the function denoted $\alpha(r)$ is half the mean of the sum of $M$ capped-binomial coordinates, and the function denoted $\mu(r)$ is $Hp(r)$. Thus their saddle equation $\alpha(r)=\mu(r)$ is precisely \eqref{ssm:eq:exactsaddle}. Our $m_0=Hp(r_M)$ lies in their interval
\[
 I_0=[H/2-M^{3/2+2\eps},\,kM/2]
\]
for any sufficiently small fixed $\eps>0$ and large $M$: its distance below $H/2$ is $\tfrac12\ell M^{3/2}+o(M^{3/2})$, and its distance below $kM/2$ has the same positive leading order. Hence the stationary-saddle expansion in their equation (29) applies. For bounded $\theta_M$, and consequently bounded $T_M$, it gives
\begin{equation}\label{ssm:eq:sourceLrate}
 L_M=L_0(T_M)+O(M^{-1}).
\end{equation}
This uses an actual quantitative source expansion, not a rate inferred from the central limit theorem.

Their Lemma 2, in this same parametrization, gives
\begin{equation}\label{ssm:eq:sourceSrate}
 S_M=\Phi(L_M+T_M)+O(M^{-1})
\end{equation}
uniformly for the bounded parameters at issue. In particular there is no separate $M^{-1/2}$ term beyond the shift in the argument. The $-1$ in $k=M/2-1+(T_M/2)\sqrt M$ is already part of that lattice-sensitive source formula.

For completeness, differentiating \eqref{ssm:eq:Lzero} yields the derivative used in \cite[equation (14)]{mww}. Write $\xi=2L_0(T)^2+L_0(T)T$. Since $\lambda'(x)=-\lambda(x)(x+\lambda(x))$, the equation $L_0(T)=\lambda(L_0(T)+T)$ implies
\[
 L_0'(T)=-\xi\{1+L_0'(T)\},\qquad
 L_0'(0)=-\frac{2\ell^2}{1+2\ell^2}.
\]
Equations \eqref{ssm:eq:sourceLrate} and \eqref{ssm:eq:sourceSrate} now give
\[
 L_M+T_M=\ell+\frac{T_M}{1+2\ell^2}+O(M^{-1}),
\]
and $\phi(\ell)/\Phi(\ell)=\ell$ proves
\begin{equation}\label{ssm:eq:Srate}
 \log S_M=\log\Phi(\ell)+\frac{2\theta_M\ell}{1+2\ell^2}M^{-1/2}
 +O(M^{-1}).
\end{equation}
Multiplying its remainder by $f\le C\sqrt M$ gives $O(M^{-1/2})=o(1)$. Substitution into \eqref{ssm:eq:criticalexact} proves \eqref{ssm:eq:criticalexplicit}, and also its $f=o(\sqrt M)$ corollary. The full enumeration still has only a qualitative $o(1)$ relative error; the auxiliary estimate \eqref{ssm:eq:Srate} does not strengthen that claim.

\subsection{The established all constrained count}
The required special case of \cite[Corollary 2]{mww} is
\begin{equation}\label{ssm:eq:allcapped}
 B(M,k)=2^H Z_0q^M
 \exp\{2\ell\theta_M\sqrt M+J\theta_M^2+o(1)\},
 \qquad \theta_M=O(1).
\end{equation}
The source actually permits $\theta_M=o(M^{1/4})$ and includes a cubic term
\[
 \frac{4\ell(6\ell^2-1)}{3(1+2\ell^2)^3}\theta_M^3M^{-1/2},
\]
which is $o(1)$ in the bounded-offset case used here. Its exponential base is written there as $e^{-\ell^2}/(\ell\sqrt{2\pi})$; equation \eqref{ssm:eq:ell} makes this exactly $q$. The constants $q,Z_0,J$ and the square-root boundary effect in \eqref{ssm:eq:allcapped} are therefore attributed to that source.

For the matrix bridge, $M=n+1$ and $k=\lfloor n/2\rfloor$. Hence
\[
 \theta_M=\begin{cases}1/2,&n\text{ even},\\0,&n\text{ odd}.
 \end{cases}
\]
Theorem~\ref{ssm:thm:transfer} at $f=1$ gives
$A_n\sim B(n+1,\lfloor n/2\rfloor)/\Phi(\ell)$. Moreover
\[
 \frac{q^{n+1}}{\Phi(\ell)}=q^ne^{-\ell^2/2},\qquad
 \sqrt{n+1}-\sqrt n=O(n^{-1/2}).
\]
Substitution into \eqref{ssm:eq:allcapped} proves Theorem~\ref{ssm:thm:matrix} with exactly the factor shown in \eqref{ssm:eq:matrix}.

The free-vertex critical correction and the matrix parity effect have related origins but different roles. The former comes from a bounded nonzero mean shift in a Gaussian quadratic expectation; the latter is inherited from the lattice-sensitive all-capped count. For one free vertex the critical corrections tend to one, but the matrix's even-order factor $e^{\ell\sqrt n}$ remains.
